% Baccalauréat malgache, série D, 2014 : énoncé.
\TCDBacAttribution{D}{2014}
\providecommand{\cm}{~\mathrm{cm}}

\noindent \underline{\textsc{Exercice I}} (5 points) \newline
      Soit le polyn\^ome $P$ \`a variable complexe $z$ d\'efini par 
      $P(z) = z^3 - (5+i) z^2 + (a + ib) z - 8 - 16 i$ o\`u $a$ et $b $ sont des nombres r\'eels non nuls.
      \begin{enumerate}
            \item D\'eterminer les valeurs de $a$ et $b$ pour que $2i$ soit une solution de l'\'equation $P(z) = 0$.
            \item R\'esoudre dans $\mathbb{C}$ l'\'equation 
                  $z^3 - (5+i) z^2 + (10 + 6i) z - 8 - 16 i = 0$.
            \item Dans le plan complexe muni d'un rep\`ere orthonorm\'e $(O ; \vec{u} , \vec{v} )$ d'unit\'e $1\cm$,
                  on consid\`ere les points $I$, $J$, $K$ et $A$ d'affixes respectives :
                  $z_I = 2i$, $z_J = 3 + i$, $z_K = 2 - 2i$ et $z_A = \sqrt{3} + i$.
                  \begin{enumerate}
                        \item Placer les points $I$, $J$ et $K$.
                              A l'aide d'un compas et d'une r\`egle non gradu\'ee, construire le point $A$.
                        \item D\'emontrer que les droites $(JI)$ et $(JK)$ sont perpendiculaires.
                  \end{enumerate}
            \item Soit $B$ le point du plan tel que $z_B = \bar{z_A}$.
                  \begin{enumerate}
                        \item Calculer les distances $OA$, $OB$ et $AB$.
                              En d\'eduire la nature du triangle $AOB$.
                        \item Calculer l'affixe du point $C$ pour que $AOBC$ soit un losange.
                  \end{enumerate}
            \item Soit $S$ la similitude plane directe qui transforme $B$ en $C$ et laisse invariant le point $O$.
                  D\'eterminer l'expression complexe de $S$ et pr\'eciser ses \'el\'ements caract\'eristiques.
      \end{enumerate}

\noindent \underline{\textsc{Exercice II}} (5 points) \newline
      Une bo\^ite contient six billets num\'erot\'es de 1 \`a 6.
      On tire au hasard, successivement et sans remise deux billets de la bo\^ite.
      On suppose que tous les tirages sont \'equiprobables.
      \begin{enumerate}
            \item Chaque r\'esultat est repr\'esent\'e par un couple $(a,b)$ de deux nombres distincts o\`u 
                  $a$ est le num\'ero apparu au premier tirage et $b$ le num\'ero apparu au deuxi\`eme tirage.
                  \begin{enumerate}
                        \item D\'emontrer qu'il y a 30 couples possibles.
                        \item Soit $A$, l'\'ev\`enement : << Les deux num\'eros tir\'es sont pairs >>.
                              Calculer la probabilit\'e de $A$.
                        \item Calculer la probabilit\'e de l'\'ev\`enement : 
                              $B$ : << Obtenir au moins un num\'ero impair >>.
                  \end{enumerate}
            \item Soit $X$ la variable al\'eatoire qui, \`a chaque r\'esultat, associe la diff\'erence entre le plus grand
                  et le plus petit des deux nombres.
                  Par exemple, pour les couples $(3,5)$ et $(5,3)$, la variable al\'eatoire $X$ prend la valeur $5 - 3 = 2$.
                  \begin{enumerate}
                        \item Quelles sont les valeurs possibles de la variable al\'eatoire $X$\,?
                        \item Calculer les probabilit\'es : $P(X = 1)$ et $P(X = 3)$.
                        \item Donner la loi de probabilit\'e de la variable al\'eatoire $X$.
                        \item Calculer l'esp\'erance math\'ematique et la variance de la variable al\'eatoire $X$.
                  \end{enumerate}
      \end{enumerate}

%%%%%%%%%%%%%%%%%%%%%%%%%%%%%%%%%%%%%%%%%%%%%%%%%%%%%%%%%%%%%%%%%%%%%%%%%%%%%%%%%%%%%%%%%%%%%%%
%%%%%%%%%%%%%%%%%%%                   %%%%%%%%%%%%%%%%%%%                   %%%%%%%%%%%%%%%%%%%
\providecommand{\Df}{D_f}
%%%%%%%%%%%%%%%%%%%%%%%%%%%%%%%%%%%%%%%%%%%%%%%%%%%%%%%%%%%%%%%%%%%%%%%%%%%%%%%%%%%%%%%%%%%%%%%
\noindent \underline{\textsc{Probleme}} \newline
      On consid\`ere la fonction num\'erique $f$ d'une variable r\'eelle r\'eelle $x$ d\'efinie sur $\Df = ] 0 ; +\infty[$
      par $f(x) = \frac{1}{2} (\ln x)^2 + \frac{2}{x}$.
      On d\'esigne par $(\mathscr{C})$ la courbe repr\'esentative de $f$ dans le rep\`ere orthonorm\'e direct 
      $(O ; \vec{\imath} , \vec{\jmath} )$ d'unit\'e $1\cm$.
      \begin{enumerate}
            \item \begin{itemize}
                        \item Calculer les limites de $f$ aux bornes de son ensemble de d\'efinition.
                        \item Calculer $\lim_{x\rightarrow+\infty} \frac{(\ln x)^2}{x}$ (poser $x = X^2$).
                              En d\'eduire que $(\mathscr{C})$ 
                              admet une branche parabolique suivant $(Ox)$ au voisinage de $+\infty$.
                  \end{itemize}
            \item On consid\`ere une deuxi\`eme fonction num\'erique $g$ d\'efinie sur $]0 ; +\infty[$ par $g(x) = x \ln x - 2$.
                  \begin{enumerate}
                        \item Calculer $\lim_{x\rightarrow 0^+ } g(x)$ et $\lim_{x\rightarrow +\infty} g(x)$.
                        \item Etudier la variation de $g$,
                              puis dresser son tableau de variation.
                        \item D\'emontrer qu'il existe un nombre r\'eel unique $\alpha \in ]2 ; e[$ tel que $g(\alpha) = 0$.
                              En d\'eduire le signe de $g(x)$. 
                  \end{enumerate}
            \item \begin{enumerate}
                        \item D\'emontrer que pour tout $x \in ] 0 ; +\infty [$, $f'(x) = \frac{g(x)}{x^2}$.
                        \item Dresser le tableau de variation de $f$.
                              On prend $\alpha = 2,5$ et $f(\alpha) = 1,12$.
                  \end{enumerate}
            \item \begin{enumerate}
                        \item D\'eterminer une \'equation de la tangente $(T)$ \`a $(\mathscr{C})$ au point d'abscisse $x_0 = 1$.
                        \item Tracer dans le m\^eme rep\`ere la droite $(T)$ et la courbe $(\mathscr{C})$.
                  \end{enumerate}
            \item \begin{enumerate}
                        \item On donne $\int_{\alpha}^{e} \ln x \dx = -\alpha \ln \alpha + \alpha$.
                              A l'aide d'une int\'egration par parties, d\'emontrer que 
                              $\int_{\alpha}^{e} \frac{1}{2} (\ln x)^2 \dx = \frac{e}{2} - \frac{\alpha}{2} (\ln \alpha)^2 
                              + \alpha \ln \alpha - \alpha$.
                        \item Calculer, en fonction de $\alpha$ et en $\cm^2$, l'aire g\'eom\'etrique $\mathscr{A}$ 
                              du domaine plan d\'elimit\'e par la courbe $(\mathscr{C})$, l'axe des abscisses et les droites
                              d'\'equations : $x = \alpha$ et $x = e$.
                        \item Sachant que : $\alpha \ln \alpha = 2$,
                              d\'emontrer que 
                              $\mathscr{A} = \left( \frac{e}{2} + 4 - \frac{6}{\alpha} - \alpha \right) \cm^2$.
                  \end{enumerate}
      \end{enumerate}







